\subsection{Attacks Taxonomy}\label{sec:cryptographic_security_overview:attacks}

Below, we list the most common categories of attacks within the context of applied cryptography.

\paragraph*{\faCog~Misconfiguration, incorrect usage, and misuse attacks.} This category of attacks exploits weaknesses arising from the improper configuration or use of algorithms rather than flaws inherent in the algorithms themselves or their implementation. This category of attacks includes secret data reuse (e.g., using the same key for multiple purposes or reusing initialization vectors or nonces), low entropy in the generation of (supposedly) random values,
%TODO (\Cref{future_sec:randomness}), 
use of weak keys (\Cref{sec:cryptography_basics:keys}), and improper key management (\Cref{sec:key_management}).

\paragraph*{\faBug~Incorrect implementation attacks.} This category of attacks exploits vulnerabilities --- due to bugs in the code or non-adherence to the original theoretical specification --- in the implementation of algorithms.

\exampleBlock{A real-world example of an incorrect implementation attack}{An example of such a category of attacks is the key reinstallation attack on WPA2 \cite{vanhoef_key_2017}.}

\paragraph*{\faShippingFast~Supply chain attacks.} This category of attacks aims to introduce vulnerabilities in hardware, software, or third-party components of cryptographic modules before their deployment. These attacks typically involve introducing malicious (and often subtle) modifications to the design and manufacture of hardware and circuits or tampering with firmware or software (e.g., open-source libraries) and their distribution to introduce cryptographic backdoors (\Cref{sec:programming_languages_and_software:trusted_computing_base}).

\remarkBlock{Supply chain attacks to cryptography algorithms}{The practice of inserting cryptographic backdoors in (usually asymmetric) algorithms or implementations while preserving their functional correctness is called \textbf{kleptography} \cite{young_dark_1996}.}

\paragraph*{\faUserSecret~Human factor attacks.} This category of attacks targets individuals rather than weaknesses in algorithms or hardware. This category of attacks includes social engineering (e.g., phishing), coercive attacks (e.g., rubber-hose cryptanalysis), malicious insiders (e.g., evil maid attacks), and also insecure workarounds driven by poor usability.

\paragraph*{\faAtom~Quantum attacks.} This category of attacks relies on quantum computers to efficiently solve mathematical problems that are assumed to be hard for classical computers.
%TODO (\Cref{future_sec:post_quantum_cryptography}).

\paragraph*{\faQuestion~Guessing attacks.} This category of attacks includes approaches based on (sometimes educated) guesses and exhaustive trial-and-error methods. This category of attacks includes \textbf{brute force attacks} (also called \textit{exhaustive key search}) --- systematically attempting all possible secret or private keys until the correct ones are discovered --- \textbf{dictionary attacks} --- systematically attempting a precompiled list of words, phrases, or common passwords (exploiting users' tendencies to choose easily guessable passwords) until the correct ones are discovered --- and \textbf{rainbow table attacks} --- relying on pre-computed tables caching the outputs of a specific \gls{chf} used for conducting pre-image attacks (\Cref{sec:hash_functions:cryptographic_hash_function}).

\warningBlock{Guessing attacks are relevant}{Although guessing attacks may seem trivial, they are a serious threat and should not be underestimated --- especially given that the computational power available to attackers keep increasing (see \textbf{Moore's law}). Intuitively, the higher the computational power, the faster a guessing attack is. Also, dedicated techniques for improving over random guessing (that is, doing educated guessing) are often available for specific ciphers to reduce the key search space and further accelerate guessing attacks. Remember that cryptography is ultimately based on one-way and trapdoor functions which are \textbf{computationally} --- \textbf{not theoretically} --- hard to invert (\Cref{sec:trapdoor_functions}). As a consequence and to mitigate the risk of guessing attacks, standardization bodies such as \gls{nist} periodically recommend to increase the length of private and secret keys to enlarge the key search space.}

\paragraph*{\faCalculator~Mathematical and algorithmic attacks.} This category of attacks aims to break the foundations of algorithms by solving the underlying hard mathematical problems (\Cref{sec:mathematical_underpinnings}) or trying to find weaknesses through cryptanalysis. This category of attacks includes factorization and discrete-logarithm solvers, differential cryptanalysis --- that is, analyzing patterns in outputs for different inputs --- and linear cryptanalysis --- that is, attempting to find linear approximations of the relationship among plaintexts, ciphertexts, and keys (\Cref{sec:cryptographic_security_overview:attack_model}).

\paragraph*{\faHourglassHalf~Side-channel attacks.} This category of attacks aims to infer secret data such as private keys from side effects caused by the execution of algorithms. This category of attacks includes power analysis attacks and timing attacks (\Cref{sec:security_of_cryptographic_algorithms:attacks:sca}).

\paragraph*{\faExclamation~Fault injection attacks.} This category of attacks consists of deliberately introducing faults into a cryptographic module to cause it to behave in unexpected ways, and consequently leak secret data. This category of attacks includes glitching attack (\Cref{sec:security_of_cryptographic_algorithms:attacks:fia}).

\paragraph*{\faDatabase~Data remanence attacks.} This category of attacks aims to recover residual physical traces of secret data that remain on a storage medium after use. This category of attacks includes cold boot attacks and forensics recovery (\Cref{sec:security_of_cryptographic_algorithms:attacks:dra}).

\paragraph*{\faEnvelope~Protocol attacks.} This category of attacks targets the composition and use of algorithms for providing security properties such as confidentiality, integrity, and authenticity (\Cref{sec:security_properties}) in the context of protocols (\Cref{sec:security_of_cryptographic_protocols}).

\remarkBlock{White-box cryptography}{Besides the aforementioned categories, a quite recent category of attacks is the so-called \textbf{white-box cryptography}, an extreme attack scenario in which the attacker has full unrestricted access to an algorithm implementation. This category of attacks was first proposed in 2002 \cite{chow_white_2002} and has many applications, including \gls{drm} and protection of cryptographic keys in the presence of malware. However, due to the huge power given to the attacker in white-box cryptography, there have been little to no practical results yet.}
